%==============================================================================
\chapter{نطاق المشروع}
%==============================================================================

\section{المجال المُغطّى}
يغطّي نطاق كراس الشروط هذا \textbf{الجزء الأول} من النظام: الإدارة اليدوية
للعمليات وتخطيطها ومتابعتها، إضافةً إلى تهيئة البنية لدمج المؤشّرات الآلية
مستقبلًا.

\section{النموذج المهني وقواعد التسيير}
\begin{xltabular}{\textwidth}{>{\bfseries}p{3cm} X}
\toprule
\textbf{الكيان} & \textbf{قواعد التسيير} \\
\midrule
\endhead
مربّي النحل & يمكن أن يملك عدة مزارع. \\
المزرعة & يمكن أن تضمّ عدة مواقع لتربية النحل. \\
الموقع & يحتوي على خلية واحدة أو أكثر، وله موقع جغرافي (خط العرض، خط الطول، الارتفاع، تاريخ التشغيل، وتواريخ محتملة للنقل / الإغلاق). يمكن لموقع أن يستقبل خلايا من مربّين ومزارع مختلفة. \\
الخلية & تتكوّن من مقصورة قاعدية (قاعدة / جسم) ومن 5 طوابق تمديد على الأكثر (عاسلات). \\
الجسم / العاسلة & يمكن أن يستقبل إطارات شمعية مركّبة على مواضع (\emph{slots}) محدّدة. يجب أن يحتوي الجسم القاعدي على إطار واحد على الأقل، والسعة القصوى للجسم كما للعاسلة هي 10 إطارات. \\
الإطار & يشغل موضعًا (\emph{slot}) وحيدًا في الجسم أو العاسلة. \\
عون تربية النحل & يقوم بالزيارات وعمليات التفتيش. يمكن متابعة خلية على التوالي من قِبَل عدة أعوان، لكنها تكون تحت مسؤولية عون واحد في لحظة معيّنة. \\
العون المشرف & يصادق على برنامج زيارة الخلايا التي هي في عهدته. \\
\bottomrule
\end{xltabular}

\begin{encadre}[nobreak=true]
\textbf{قيد التركيب:} خلية واحدة = جسم/قاعدة إجباري واحد + من 0 إلى 5 عاسلات،
كل جسم/عاسلة = من 1 إلى 10 إطارات، إطار واحد لكل موضع.
\end{encadre}

\begin{figure}[htbp]
\centering
\begin{tikzpicture}[
  every node/.style={font=\small},
  ent/.style={draw=mielfonce,fill=grisdoux,rounded corners=3pt,minimum height=9mm,inner xsep=7pt,thick},
  fl/.style={-{Latex[length=2mm]},mielfonce,thick}
]
\node[ent] (a) {مربّي النحل};
\node[ent,right=7mm of a] (b) {المزرعة};
\node[ent,right=7mm of b] (c) {الموقع};
\node[ent,right=7mm of c] (d) {الخلية};
\node[ent,below=9mm of d] (e) {الجسم / العاسلة};
\node[ent,left=7mm of e] (f) {الإطار};
\draw[fl] (a)--(b); \draw[fl] (b)--(c); \draw[fl] (c)--(d);
\draw[fl] (d)--(e); \draw[fl] (e)--(f);
\node[font=\scriptsize,mielfonce,above=0.5mm] at ($(a)!0.5!(b)$) {1..n};
\node[font=\scriptsize,mielfonce,above=0.5mm] at ($(b)!0.5!(c)$) {1..n};
\node[font=\scriptsize,mielfonce,above=0.5mm] at ($(c)!0.5!(d)$) {1..n};
\end{tikzpicture}
\caption{التسلسل الهرمي للكيانات المهنية لنظام Zümm.}
\label{fig:hierarchie}
\end{figure}

\begin{figure}[htbp]
\centering
\begin{tikzpicture}[font=\small]
\draw[draw=mielfonce,thick,rounded corners=6pt,fill=grisdoux] (0,0) rectangle (11,4.4);
\node[anchor=north east,mielfonce,font=\bfseries] at (10.8,4.25) {موقع لتربية النحل};
\node[anchor=north east,font=\footnotesize] at (10.8,3.7)
  {خط العرض / خط الطول / الارتفاع، تاريخ التشغيل};
\foreach \x in {0.9,4.3,7.7}{
  \begin{scope}[shift={(\x,0.7)}]
    \fill[miel!30,draw=mielfonce,thick] (0,0) rectangle (1.4,0.6);
    \fill[miel!40,draw=mielfonce,thick] (0,0.6) rectangle (1.4,1.2);
    \fill[mielfonce] (-0.12,1.2)--(1.52,1.2)--(0.7,1.7)--cycle;
    \draw[mielfonce,thick] (-0.1,-0.12) rectangle (1.5,0);
  \end{scope}
}
\node[font=\scriptsize] at (1.6,0.45) {خلية};
\node[font=\scriptsize] at (5.0,0.45) {خلية};
\node[font=\scriptsize] at (8.4,0.45) {خلية};
\node[anchor=west,font=\footnotesize,align=right] at (11.4,3.2)
  {موارد\\(رحيق / غبار طلع)\\على بُعد 1 كم تقريبًا};
\draw[-{Latex[length=2mm]},mielfonce,dashed] (11,3.0) -- (11.35,3.0);
\node[anchor=west,font=\footnotesize,align=right] at (11.4,1.2)
  {مسافات الأمان\\100 / 200 / 300 م\\3 إلى 4 كم عن\\المزروعات المعالَجة};
\end{tikzpicture}
\caption{مخطّط موقع لتربية النحل: عدة خلايا محدّدة الموقع وقيود التموضع.}
\label{fig:site}
\end{figure}

\section{المستخدمون وحقوق الوصول}
\begin{xltabular}{\textwidth}{>{\bfseries}p{3.2cm} X}
\toprule
\textbf{الملف} & \textbf{الدور والحقوق} \\
\midrule
\endhead
مربّي النحل / المالك & يطّلع على لوحات قيادة الإنتاج والمردودية والخلاصة، ويقرّر إغلاق موقع أو نقل الخلايا. \\
عون تربية النحل & يقوم بالزيارات ويملأ التقارير وينفّذ العمليات على الخلايا التي في مسؤوليته. \\
العون المشرف & يصادق على برامج زيارة الخلايا التي هو مسؤول عنها. \\
المسؤول / المدير & يراقب التنبيهات (انخفاض الإنتاج أو العدد) ويُطلق عمليات التفتيش الوشيكة. \\
مدير النظام & يدير معطيات النظام ووحدات القياس والعتبات والمرجعيات. \\
\bottomrule
\end{xltabular}

\section{خارج النطاق}
\begin{itemize}
  \item الاقتناء المادي والنشر العتادي للمستشعرات (مشروع لاحق).
  \item الأتمتة الآنية لجمع المؤشّرات (واجهة مُقرّرة لكن غير مُنفّذة في هذه المرحلة).
\end{itemize}

\section{آفاق التطوّر}
خارج نطاق هذه المرحلة، يمكن دراسة التطوّر التالي:
\begin{itemize}
  \item \textbf{تطبيق جوّال أصيل منفصل (\textenglish{React Native})، متميّز عن تطبيق
  الويب التقدّمي (\textenglish{PWA}).} الحل المعتمَد هو تطبيق ويب تقدّمي
  (\textenglish{PWA}) يضمن مسبقًا التكافؤ بين الويب والجوّال والعمل دون اتصال (انظر
  ملحق التقنيات). لن يُشرَع في تطبيق أصيل منفصل إلا إذا فرضته حاجة خاصة بالأصيل:
  إشعارات فورية موثوقة على \textenglish{iOS}، أو تحديد الموقع في الخلفية، أو التوزيع
  عبر متاجر التطبيقات. سيُعيد استخدام النواة المنطقية بلغة \textenglish{TypeScript}
  وواجهة \textenglish{REST}، لكنه سيُشكّل قاعدة شيفرة واجهة ثانية تجب صيانتها --- لا
   يُشرَع فيه إلا بطلب صريح من العميل.
  \item \textbf{وحدة الإنتاج والحصاد.}
  تتبّع الحصاد (الوزن، العائد، التصنيف) جوهري في إدارة المنحلات وهو غائب عن
  النطاق الحالي. تسمح هذه الوحدة بتسجيل حصد كل خلية ولفصل، وحساب العوائد،
  ومراقبة جودة العسل (الغبار، الرطوبة، الموصلية) وإعداد تقارير الإنتاج
  للدعم والتفقّد الصحي.

  \item \textbf{تتبّع أصل العسل --- التوجيهة الأوروبية لوجبة الفطور (2026).}
  تفرض التوجيهة الأوروبية تتبّع أصل العسل بالكامل بحلول يونيو 2026: ترقيم
  الدفعات، سجل المعالجات الصحية، الشهادات (عضوي، مؤشر جغرافي محمي، اسم
  مغربي محمي)، رمز \textenglish{QR} على البرطمان مرتبط بسجل الخلية. لا تقدم
  أي منصة منافسة وحدة تتبّع متكاملة، مما يمثّل ميزة تنافسية كبرى لـ\textenglish{Zümm}.

  \item \textbf{دمج بيانات الأرصاد الجوية.}
  تؤثّر الحرارة والرطوبة والرياح والأشعة الشمسية بشكل مباشر على سلوك المستعمرات
  (التنشّب، جمع الرحيق، استهلاك المخزون). يتيح دمج بيانات الطقس مع
  مستشعرات الخلايا إنذارات سياقية وتحسين تخطيط الزيارات.

  \item \textbf{تصور خرائطي عبر \textenglish{PostGIS}.}
  تعتمد الخلفية على \textenglish{PostgreSQL/PostGIS}، وهي قدرة غير موجودة لدى
  المنافسين. استغلال قاعدة البيانات المكانية لخريطة تفاعلية للمنحلات (التجمع،
  المسافات بين الخلايا، تحسين مسارات الزيارات، خرائط الأزهار عبر الأقمار
  الصناعية) يمثّل ميزة تنافسية كبيرة.

  \item \textbf{موصّل إنترنت الأشياء المعياري.}
  بدلاً من تطوير العتاد، إنشاء نقطة استيعاب معيارية (\texttt{/api/v1/telemetry})
  مع مصادقة برمز الجهاز ومخطط بيانات موحّد (وزن، حرارة، رطوبة، صوتي). يتيح ذلك
  دمج مستشعرات موجودة (\textenglish{BroodMinder, ApisProtect, BEEP, ESP32 DIY})
  دون تطوير العتاد.

  \item \textbf{الكشف الصوتي والرؤية الحاسوبية.}
  تُظهر الأبحاث الحديثة دقّة تتراوح بين 85 إلى 99\,\% للكشف الصوتي عن السوس
  والتنشّب، فضلاً عن العدّ الآلي للسبّورة اللزجة بالرؤية
  (\textenglish{YOLOv11 Nano}، $R^2 = 0{,}99$). يمكن دمج هذه النماذج مفتوحة
  المصدر في وحدة تفقّد الصور ومحلّل صوتي، مما يُثري تقارير الزيارات بمؤشرات
  آلية.

  \item \textbf{التوأم الرقمي للمستعمرة وائتمانات الكربون.}
  على المدى الأبعد، تتيح النمذجة التنبؤية (التوائم الرقمية) محاكاة
  \emph{ماذا لو} حول معالجة السوس أو التغذية أو الهجرة. يفتح تتبّع
  التلقيح الدقيق وحساب البصمة الكربونية للنحل الباب أمام أسواق ائتمانات
  الكربون والدعم البيئي الأوروبي.
\end{itemize}
